\documentclass{rvcepaper}
\begin{document}

% ---------- HS351TA Entrepreneurship & IPR (Improvement CIE)
\begin{iempaperB}
\paperinfo{ayear={2024-2025 (ODD Sem)},date={27\textsuperscript{th} January 2025},code={HS351TA},sem={V Semester},exam={Improvement CIE},maxmarks={10 + 50},duration={120 Mins},
 title={ENTREPRENEURSHIP AND INTELLECTUAL PROPERTY RIGHTS}}
\iemheadB
\ptitle{PART -- A}
\begin{cietable}{M}{BT}{CO}
\q{1}{\blank[3.2cm] is a fundamental framework in marketing that helps businesses identify and reach their most valuable customers.}{1}{1}{2}
\q{2}{Creating a \blank[3.2cm] is about identifying and clearly communicating what makes a product, service, or brand distinctively valuable to its target audience}{1}{1}{4}
\q{3}{A company is launching a premium product with a production cost of Rs 30 per unit. They estimate fixed costs for branding and marketing at Rs 2,00,000. The company plans to sell 10,000 units in the first year and targets a 40\% profit margin. What should be the selling price per unit to achieve the desired profit margin?}{2}{2}{2}
\q{4}{Toyota developed the \textbf{Toyota Way}, a set of principles that focuses on \blank[3.2cm] and respect for people.}{1}{1}{4}
\q{5}{Differentiate between Planning and Strategy.}{2}{1}{2}
\q{6}{With \blank[3cm] growth, a company expands through its own operations using its own internal resources.}{1}{1}{4}
\q{7}{If a company's gross profit margin is increasing but its net profit margin is declining, what could be the reason?}{2}{2}{3}
\end{cietable}
\ptitle{PART -- B}
\begin{cietable}{M}{BT}{CO}
\q{1.}{A tech company is developing a mobile app that combines fitness tracking with AI-based personalized coaching. They are seeking funding from venture capitalists. How can the company effectively articulate its unique value proposition in the business plan to stand out from competitors in the crowded fitness tech market?}{10}{4}{4}
\q{2.}{a)\ A start-up has developed a patent-pending water purification device that uses renewable energy. They aim to market it in developing countries. How should they address intellectual property, innovation, and scalability in their business plan to appeal to investors?\newline
b)\ A global e-commerce platform is considering a strategic alliance with a logistics company to enhance delivery efficiency in underserved regions. The e-commerce platform wants to improve delivery times, and the logistics company seeks access to a broader customer base. What factors should both companies consider when forming a strategic alliance, and how can they ensure that the partnership is mutually beneficial?}{5\newline\newline\newline\newline\newline\newline 5}{4\newline\newline\newline\newline\newline\newline 3}{4\newline\newline\newline\newline\newline\newline 4}
\q{3.}{What are the objectives of Financial Management? Investigate Financial Management Tools and Techniques.}{10}{2}{2}
\q{4.}{Illustrate how successful companies have leveraged each of the 4Ps to achieve market success.}{10}{3}{2}
\q{5.}{A manufacturing company needs Rs 20 Lacs to purchase advanced machinery. They are considering two options:
\begin{bul}
\item Taking a bank loan with a 7\% annual interest rate.
\item Selling 20\% of the company's equity to a private investor.
\end{bul}
\begin{num}
\item What factors should the company consider when choosing between equity and debt financing? Which option might be better if the company expects rapid growth? Justify with a case study.
\item What steps can the company take to manage debt repayment and avoid defaulting on the loan?
\item Illustrate how successful companies have leveraged each of the 4Ps to achieve market success.
\end{num}}{10}{4}{2}
\end{cietable}
\btfoot{|c|c|c|*{11}{C{0.82cm}|}}{\hline
\multirow{3}{*}{\bfseries\shortstack{Marks Distribution}} & \multicolumn{2}{c|}{\bfseries Particulars} & CO1 & CO2 & CO3 & CO4 & CO5 & L1 & L2 & L3 & L4 & L5 & L6\\\cline{2-14}
 & \multirow{2}{*}{Max Marks} & Quiz & -- & 5 & 2 & 3 & -- & 06 & 04 & -- & -- & -- & --\\\cline{3-14}
 & & Test & -- & 20 & 10 & 20 & -- & -- & 10 & 15 & 25 & -- & --\\\hline}
\stars{*****}
\end{iempaperB}

% ---------- IM352IA Operations Management (CIE - III)
\begin{iempaperB}
\paperinfo{ayear={2024-2025 (ODD Sem)},date={27/01/2025},code={IM352IA},sem={V},exam={CIE -- III},maxmarks={10 + 50},duration={30 + 90 Min},
 title={OPERATIONS MANAGEMENT}}
\iemheadB
\begin{cietable}{M}{BT}{CO}
\qpart{Part -- A}
\q{1}{A \blank[2cm] is any resource whose capacity is less than or equal to the demand placed on it.}{1}{1}{3}
\q{2}{In product mix decisions, TOC prioritizes products based on \blank[2cm] per unit of the bottleneck resource.}{1}{1}{3}
\q{3}{The Theory of Constraints uses the \blank[2cm] process to repeat improvements once a bottleneck is resolved.}{1}{1}{3}
\q{4}{In manufacturing, reducing \blank[2cm] time at a bottleneck can significantly improve efficiency.}{1}{1}{3}
\q{5}{Differentiate between Static Schedule and Dynamic Schedule}{2}{2}{4}
\q{6}{What is bottleneck scheduling, and why is it critical?}{2}{2}{4}
\q{7}{Using the Earliest Due Date (EDD) rule, determine the sequence for the following jobs based on their due dates:
\qtabc{|c|c|c|}{\hline Job & Processing time(hr) & Due date(hr)\\\hline A & 4 & 10\\\hline B & 3 & 8\\\hline C & 2 & 6\\\hline D & 5 & 12\\\hline}}{2}{2}{4}
\qpart{Part --B}
\q{1}{Explain how the Theory of Constraints can be applied to manage bottlenecks in a manufacturing environment. Provide examples of how identifying and addressing bottlenecks can improve overall production efficiency.}{10}{3}{3}
\q{2}{Consider the following flow shop scheduling problem with two machines (Machine 1 and Machine 2) and five jobs. The processing times (in hours) for each job on the two machines are as follows:
\qtabc{|c|c|c|}{\hline Job & Machine 1 & Machine 2\\\hline A & 5 & 4\\\hline B & 7 & 6\\\hline C & 3 & 5\\\hline D & 4 & 3\\\hline E & 6 & 7\\\hline}
Using Johnson's Rule, find the optimal job sequence that minimizes the makespan.}{10}{3}{4}
\q{3}{Consider the following 3 machines and 5 jobs flow shop problem.
\qtabc{|c|c|c|c|}{\hline Job, $i$ & \multicolumn{3}{c|}{Processing Time}\\\hline & Machine 1 & Machine 2 & Machine 3\\\hline 1 & 10 & 10 & 12\\\hline 2 & 12 & 8 & 20\\\hline 3 & 16 & 6 & 14\\\hline 4 & 12 & 4 & 20\\\hline 5 & 20 & 8 & 8\\\hline}
Check whether the Johnson's algorithm can be extended to this problem. If yes, find the optimal sequence and the corresponding makespan.}{15}{3}{4}
\q{4}{Use graphical method to minimize the time needed to process the following jobs on the machines shown (i.e. for each machine, find the job which should be scheduled first). Also, calculate the total time elapsed to complete both jobs.
\qtabc{|l|l|c|c|c|c|c|}{\hline Job 1 & Sequence & A & B & C & D & E\\\cline{2-7} & Time (Hrs). & 4 & 3 & 6 & 2 & 7\\\hline Job 2 & Sequence & C & B & E & D & A\\\cline{2-7} & Time (Hrs) & 6 & 3 & 5 & 3 & 7\\\hline}}{15}{4}{4}
\end{cietable}
\btfoot{|c|c|c|*{10}{C{0.85cm}|}}{\hline
\multirow{3}{*}{\bfseries Marks Distribution} & \multicolumn{2}{c|}{\bfseries Particulars} & CO1 & CO2 & CO3 & CO4 & L1 & L2 & L3 & L4 & L5 & L6\\\cline{2-13}
 & \multirow{2}{*}{Max Marks} & Quiz & -- & -- & 4 & 6 & 4 & 6 & -- & -- & -- & --\\\cline{3-13}
 & & Test & -- & -- & 10 & 40 & -- & -- & 35 & 15 & -- & --\\\hline}
\stars{*******}
\end{iempaperB}

% ---------- IM353IA Quality Assurance (Improvement CIE)
\begin{iempaperB}
\paperinfo{ayear={2024-2025 (ODD Sem)},usn=0,date={28\textsuperscript{th} January 2025},code={IM353IA},sem={V},exam={Improvement CIE},maxmarks={50},duration={90 Min},
 title={QUALITY ASSURANCE}}
\iemheadB
\ptitlem{PART -- A}{Max Marks: 10}
\begin{cietable}{M}{BT}{CO}
\q{1}{Three components each with a reliability of 0.9 are placed in series. What is the reliability of the system?}{2}{L2}{CO1}
\q{2}{There are two components that are identical in a parallel configuration where the second unit is in standby and will be activated only upon failure of the primary unit. The switching to the backup unit is considered certain. The failure rate for each component is 0.002 failures/hour. What is the reliability for 200 hours?}{2}{L1}{CO2}
\q{3}{The failure rates for two units in a parallel configuration, where the secondary unit is on back up until switched on upon failure of the primary unit, are each 0.002 failures/hour and the switch reliability is 0.95 per operation. What is the reliability for 300 hours?}{2}{L2}{CO3}
\q{4}{List the Advantages of CRD}{2}{L2}{CO3}
\q{5}{Expand ANOVA and ANCOVA}{2}{L2}{CO2}
\end{cietable}
\ptitlem{PART -- B}{Max Marks: 50}
\begin{cietable}{M}{CO}{BT}
\q{1}{Define the following terms: Reliability; MTTF; MTBF; MTTR; Availability ,Maintainability}{10}{CO3}{3}
\q{2}{Explain the utility of `bath tub curve'. What are the reasons for failure during the three stages of equipment life?}{10}{CO1}{3}
\q{3}{Consider the following series-parallel configuration:
\par\vspace{4pt}\centerline{\begin{tikzpicture}[x=1cm,y=1cm,font=\small,every node/.style={draw,minimum width=0.45cm,minimum height=0.45cm,inner sep=1pt}]
\node (a) at (1,1.4) {R}; \node (b) at (3.0,2.1) {R}; \node (c) at (3.0,0.7) {R}; \node (d) at (4.7,1.4) {R};
\node (e) at (1,-0.7) {R}; \node (f) at (3.0,-0.7) {R}; \node (g) at (4.7,-0.7) {R}; \node (h) at (7.2,0.35) {R};
\draw[-{Stealth}] (-0.4,0.35) -- (0,0.35); \draw (0,-0.7) -- (0,1.4);
\draw (0,1.4) -- (a) -- (2,1.4); \draw (2,0.7) -- (2,2.1); \draw (2,2.1) -- (b) -- (4,2.1); \draw (2,0.7) -- (c) -- (4,0.7); \draw (4,0.7) -- (4,2.1); \draw (2,1.4) -- (4,1.4) -- (d) -- (5.7,1.4);
\draw (0,-0.7) -- (e) -- (f) -- (g) -- (5.7,-0.7); \draw (5.7,-0.7) -- (5.7,1.4); \draw[-{Stealth}] (5.7,0.35) -- (h) -- (8.3,0.35);
\end{tikzpicture}}\par\vspace{4pt}
Compute the System reliability for $R=0.96$.}{10}{CO2}{3}
\q{4}{List the General practical steps and guidelines for planning and conducting DOE}{10}{CO4}{4}
\q{5.a}{How do the principles of replication and randomization work together to strengthen the internal validity of an experiment?}{5}{CO2}{3}
\q{5.b}{What are the key assumptions behind the application of a full factorial design and how do they affect the interpretation of experimental results?}{5}{CO2}{3}
\end{cietable}
\btfoot{|c|c|c|*{9}{C{0.85cm}|}}{\hline
\multirow{3}{*}{\bfseries\shortstack{Marks\\Distribution}} & \multicolumn{2}{c|}{\bfseries Particulars} & CO1 & CO2 & CO3 & CO4 & CO5 & L1 & L2 & L3 & L4\\\cline{2-12}
 & \multirow{2}{*}{\shortstack{Max\\Marks}} & Quiz & 02 & 04 & 04 & -- & -- & 02 & 08 & -- & --\\\cline{3-12}
 & & Test & 10 & 20 & 10 & 10 & -- & -- & 30 & 10 & 10\\\hline}
\stars{*****}
\end{iempaperB}

% ---------- IM254TA Financial Accounting and Costing (Improvement Test)
\begin{iempaperB}
\paperinfo{ayear={2024-2025 (ODD Sem)},date={28th January 2025},code={IM254TA},sem={V Semester},exam={Improvement Test},maxmarks={10 + 50},duration={120 Mins},
 title={FINANCIAL ACCOUNTING AND COSTING}}
\iemheadB
\ptitle{PART -- A}
\begin{cietable}{M}{BT}{CO}
\q{1}{Mention the main purpose of a Profit and Loss Account?}{02}{2}{CO3}
\q{2}{List the two main sections of a Balance Sheet?}{02}{1}{CO3}
\q{3}{Differentiate between Cost of Production and Cost of Sales?}{02}{2}{CO1}
\q{4}{What are the main components of Batch Costing?}{02}{1}{CO4}
\q{5}{What is the meaning of `Equivalent Units' in process costing?}{02}{2}{CO4}
\end{cietable}
\ptitle{PART -- B}
\begin{cietable}{M}{BT}{CO}
\q{1.}{Prepare a Trading and Profit and Loss Account for the year ending 31\textsuperscript{st} December 2024 from the following information. Also interpret the financial status of the organization.
\qtabc{|l|r|}{\hline \multicolumn{1}{|c|}{\textbf{Particulars}} & \multicolumn{1}{c|}{\textbf{Amount (\rupee)}}\\\hline Opening Stock & 1,00,000\\\hline Purchases & 5,00,000\\\hline Sales & 8,00,000\\\hline Wages & 50,000\\\hline Carriage Inwards & 20,000\\\hline Salaries & 60,000\\\hline Office Rent & 25,000\\\hline Depreciation on Furniture (10\%) & 15,000\\\hline Interest on Loan & 10,000\\\hline Closing Stock & 1,50,000\\\hline}}{10}{4}{CO3}
\q{2.}{Prepare a Balance Sheet for ABC \& Co. as on 31\textsuperscript{st} December 2024 using the following details:
\begin{bul}
\item Capital: \rupee 3,00,000
\item Profit for the Year: \rupee 50,000
\item Machinery: \rupee 2,50,000
\item Creditors: \rupee 1,00,000
\item Cash in Hand: \rupee 20,000
\item Debtors: \rupee 1,00,000
\item Bank Loan: \rupee 50,000
\item Stock: \rupee 1,30,000
\end{bul}}{10}{3}{CO3}
\q{3.}{Prepare a Cost Sheet for XYZ Ltd. for the year ending 31\textsuperscript{st} March 2024 from the following data:
\begin{bul}
\item Raw Material Purchased: \rupee 2,00,000
\item Freight on Raw Materials: \rupee 20,000
\item Direct Labor: \rupee 1,50,000
\item Direct Expenses: \rupee 25,000
\item Factory Overheads: \rupee 1,00,000 (60\% of Direct Labor)
\item Office Overheads: \rupee 75,000 (10\% of Factory Cost)
\item Selling Overheads: \rupee 50,000
\item Distribution Overheads: \rupee 30,000
\item Profit: \rupee 1,20,000
\end{bul}}{10}{3}{CO4}
\q{4.}{XYZ Ltd. is working on Job No. 45. The following details are provided for the job:
\begin{bul}
\item Direct Materials Used: \rupee 50,000
\item Direct Wages: \rupee 30,000
\item Machine Hours Worked: 20 hours
\item Machine Hour Rate: \rupee 500 per hour
\item Office Overheads: 20\% of Prime Cost
\item Profit Markup: 25\% on Total Cost
\end{bul}
Prepare a Job Cost Sheet for Job No. 45.}{10}{3}{CO4}
\q{5.}{XYZ Ltd. uses process costing to manufacture goods that go through three processes: A, B, and C. The following information is for Process B for the month of January 2025:
\begin{bul}
\item Opening Work-in-Progress (WIP): 500 units valued at \rupee 30,000 (Material: \rupee 20,000; Labor: \rupee 5,000; Overheads: \rupee 5,000)
\item Raw Materials Added: \rupee 40,000
\item Direct Labor: \rupee 25,000
\item Factory Overheads: \rupee 12,500
\item Units Introduced from Process A: 5,000
\item Units Transferred to Process C: 4,800
\item Closing WIP: 700 units (Material 100\% complete, Conversion 50\% complete)
\end{bul}
Prepare the Cost Sheet for Process B.}{10}{3}{CO4}
\end{cietable}
\btfoot{|c|c|c|*{11}{C{0.82cm}|}}{\hline
\multirow{3}{*}{\bfseries Marks Distribution} & \multicolumn{2}{c|}{\bfseries Particulars} & CO1 & CO2 & CO3 & CO4 & CO5 & L1 & L2 & L3 & L4 & L5 & L6\\\cline{2-14}
 & \multirow{2}{*}{Max Marks} & Quiz & 02 & -- & 04 & 04 & -- & 02 & -- & 04 & 04 & -- & --\\\cline{3-14}
 & & Test & -- & -- & 20 & 30 & -- & -- & -- & -- & 40 & 10 & --\\\hline}
\stars{*****}
\end{iempaperB}

% ---------- IM355TBD Advanced Decision Modelling (CIE - Improvement)
\begin{iempaperB}
\paperinfo{ayear={2024-2025 (ODD Sem)},usn=0,date={29\textsuperscript{th} January 2025},code={IM355TBD},sem={V},exam={CIE -- Improvement},maxmarks={50},duration={90 Min},
 title={ADVANCED DECISION MODELLING (Professional Core Elective)}}
\iemheadB
\begin{cietable}{M}{BT}{CO}
\qpart{Part A}
\q{1.}{List two functional roles of inventory.}{02}{1}{1}
\q{2.}{Identify two key factors in inventory problem analysis.}{02}{2}{3}
\q{3.}{How does an EOQ model without shortages work?}{02}{2}{2}
\q{4.}{How does inventory contribute to production efficiency?}{02}{2}{2}
\q{5.}{How is Markov Analysis used in business decision-making?}{02}{2}{3}
\qpart{Part B}
\q{1.}{The production department of a company requires 3,600 kg of raw material for manufacturing a particular item per year. It has been estimated that the cost of placing an order is Rs 36 and the cost of carrying inventory is 25 per cent of the investment in the inventories. The price is Rs 10 per kg. Help the purchase manager to determine an ordering policy for raw material.}{10}{1}{1}
\q{2.}{Each unit of an item costs a company Rs 40. Annual holding costs are 18 per cent of unit cost of the item due to miscellaneous changes: 1 per cent for insurance, 2 per cent allowances for obsolescence, Rs 2 for building overheads, Rs 1.50 for damage and loss, and Rs 4 miscellaneous costs. The annual demand for the item is constant at 1,000 units. Placing each order costs, the company Rs 100.
\begin{romi}
\item Calculate EOQ and the total costs associated with stocking the item.
\item If the supplier of the item will only deliver batches of 250 units, how are the stock holding costs affected?
\item If the supplier relaxes his order size requirement, but the company has limited warehouse space and can stock a maximum of 100 units at any time, what would be the optimal ordering policy and associated costs?
\end{romi}}{10}{2}{2}
\q{3.}{A product is sold at the rate of 50 pieces per day and is manufactured at a rate of 250 pieces per day. The set-up cost of the machines is Rs 2,000 and the storage cost is found to be Re 0.15 per piece per day. With labour charges of Rs 3.20 per piece, material cost at Rs 2.10 per piece and overhead cost of Rs 4.10 per piece, find the minimum cost batch size if the interest charges are 8 per cent (assume 300 working days in a year). Compute the optimal number of cycles required in a year for manufacturing of this product.}{10}{3}{3}
\q{4.}{A commodity is to be supplied at a constant rate of 200 units per day. Supplies of any amount can be obtained at any required time, but each ordering costs Rs 50; cost of holding the commodity in inventory is Rs 2 per unit per day while the delay in the supply of the item induces a penalty of Rs 10 per unit per day. Find the optimal policy (Q, t), where t is the reorder cycle period and Q is the inventory after reorder. What would be the best policy to adopt if the penalty cost becomes infinite?}{10}{3}{3}
\q{5.}{List and explain the characteristics of Markov Chain}{10}{2}{2}
\end{cietable}
\btfoot{|c|c|c|*{10}{C{0.85cm}|}}{\hline
\multirow{3}{*}{\bfseries Marks Distribution} & \multicolumn{2}{c|}{\bfseries Particulars} & CO1 & CO2 & CO3 & CO4 & L1 & L2 & L3 & L4 & L5 & L6\\\cline{2-13}
 & Quiz & Max Marks & 02 & 04 & 04 & -- & 02 & 08 & -- & -- & -- & --\\\cline{2-13}
 & Test & Max Marks & 10 & 20 & 20 & -- & 10 & 20 & 20 & -- & -- & --\\\hline}
\stars{*******}
\end{iempaperB}

% ---------- IM255TBB Enterprise Information Systems (Improvement CIE)
\begin{iempaperB}
\paperinfo{ayear={2024-2025(Odd Sem)},usn=0,hdr=2,rule=0,dept={INDUSTRIAL ENGINEERING AND MANAGEMENT},date={29\textsuperscript{th} January 2025},code={IM255TBB},sem={V Semester},exam={Improvement CIE},maxmarks={10+50},duration={120 Min},
 title={ENTERPRISE INFORMATION SYSTEMS}}
\iemheadB
\ptitle{PART -- A}
\begin{cietable}{M}{BT}{CO}
\q{1}{Define Business Process Automation (BPA).}{02}{L2}{CO3}
\q{2}{Name two common tools used in Robotic Process Automation (RPA).}{02}{L2}{CO3}
\q{3}{Mention one advantage of implementing automation in business processes.}{02}{L1}{CO2}
\q{4}{What is a key characteristic of structured data?}{02}{L1}{CO1}
\q{5}{What are the three types of data analytics?}{02}{L2}{CO2}
\end{cietable}
\ptitle{PART -- B}
\begin{cietable}{M}{BT}{CO}
\q{1.}{A mid-sized manufacturing company faces delays in order processing due to repetitive manual tasks such as data entry and validation. Propose a solution using Robotic Process Automation (RPA). Discuss how implementing RPA can streamline the process and improve efficiency.}{10}{L3}{CO2}
\q{2.}{A retail company plans to automate its invoice processing system using RPA. Explain the steps involved in the RPA implementation process for this scenario.}{10}{L3}{CO3}
\q{3.}{Design a simplified approach to integrating Business Process Automation (BPA) in the inventory management process of an e-commerce company. Include key steps and expected outcomes.}{10}{L4}{CO4}
\q{4.}{A retail chain plans to improve its inventory management system using data warehousing. Explain how a data warehouse can help in managing inventory more effectively. Include key components of a data warehouse and their roles.}{10}{L3}{CO3}
\q{5.}{Explain how data analytics in Enterprise Information Systems (EIS) can be used to improve the decision-making process for a manufacturing company aiming to reduce operational costs.}{10}{L2}{CO2}
\end{cietable}
\btfoot{|c|l|*{10}{C{0.85cm}|}}{\hline
\multirow{3}{*}{\shortstack{Marks\\Distribution}} & \textbf{Particulars} & CO1 & CO2 & CO3 & CO4 & L1 & L2 & L3 & L4 & L5 & L6\\\cline{2-12}
 & Quiz Marks & 02 & 04 & 04 & -- & 04 & 06 & -- & -- & -- & --\\\cline{2-12}
 & Test Marks & -- & 10 & 20 & 20 & -- & 10 & 30 & 10 & -- & --\\\hline}
\end{iempaperB}

\end{document}
